% Part: first-order-logic
% Chapter: natural-deduction
% Section: proving-things

\documentclass[../../../include/open-logic-section]{subfiles}

\begin{document}

\iftag{FOL}
      {\olfileid{fol}{ntd}{pro}}
      {\olfileid{pl}{ntd}{pro}}

\olsection{د \usetoken{P}{derivation} بېلګې}

\begin{ex}
راځئ چې د !!{sentence}~$(!A \land !B) \lif !A$ يو !!a{derivation}
ورکړو۔

په پيل کښې مطلوبه پايله د !!{derivation} په لاندې برخه کښې ليکو۔
\begin{prooftree}
\AxiomC{}
\UnaryInfC{$(!A\land !B) \lif !A$}
\end{prooftree}

اوس بايد معلومه کړو چې د کوم ډول استنتاج پايله د دې بڼې
!!a{sentence} کېدے شي۔ د پايلې !!{main operator}~$\lif$ دے، نو هڅه
کوو چې د~\Intro{\lif} قاعدې په کارولو پايلې ته ورسېږو۔ غوره ده چې
اړوندې فرضيې وليکو او د پرمختګ پر مهال د استنتاج قاعدو ته نښانونه
ورکړو، څو د ثبوت په پای کښې په اسانۍ ووينو چې ټولې فرضيې
!!{discharged} شوې دي که نۀ۔
\begin{prooftree}
\AxiomC{$\Discharge{!A \land !B}{1}$}
\DeduceC{$!A$}
\DischargeRule{\Intro{\lif}}{1} 
\UnaryInfC{$(!A\land !B) \lif !A$}
\end{prooftree}

اوس بايد له فرضيې~$!A \land !B$ څخه تر~$!A$ پورې پړاوونه ډک کړو۔
ځکه چې يوازې له يوه رابط~$\land$ سره کار لرو، بايد د~$\land$ د
ايستلو قاعده وکاروو۔ له دې څخه لاندې ثبوت لاس ته راځي:
\begin{prooftree}
\AxiomC{$\Discharge{!A \land !B}{1}$}
\RightLabel{\Elim{\land}}
\UnaryInfC{$!A$}
\DischargeRule{\Intro{\lif}}{1} 
\UnaryInfC{$(!A\land !B) \lif !A$}
\end{prooftree}
اوس د~$(!A \land !B) \lif !A$ يو سم !!{derivation} لرو۔
\end{ex}

\begin{ex}
اوس راځئ چې د~$(\lnot !A \lor !B) \lif (!A \lif !B)$ يو
!!a{derivation} ورکړو۔

په پيل کښې مطلوبه پايله د !!{derivation} په لاندې برخه کښې ليکو۔
\begin{prooftree}
\AxiomC{}
\UnaryInfC{$(\lnot !A \lor !B) \lif (!A \lif !B)$}
\end{prooftree}
د هغې منطقي قاعدې د موندلو دپاره چې دا پايله راکولے شي، په پايله
کښې منطقي رابطونو~$\lnot$، $\lor$ او~$\lif$ ته ګورو۔ اوس يوازې د
$\lif$ لومړۍ پېښه راته مهمه ده، ځکه همدا د وروستۍ پايلې د
!!{sentence}~!!{main operator} دے؛ $\lnot$، $\lor$ او د~$\lif$
دويمه پېښه د بل رابط په ساحه کښې دي، نو وروسته به ئې وڅېړو۔ له دې
امله د~\Intro{\lif} له قاعدې پيل کوو۔ سمه کارونه ئې بايد داسې وي:
\textbf{د ژباړې د نسخې سمون (OLND-004):} سرچينه دلته د فطري
استنتاج وروستۍ جمله “په وروستي سېکوېنټ کښې” بولي؛ دلته “وروستۍ
پايله” راوړل شوې، ځکه دا ونه د جملو ثبوت دے او سېکوېنټ نۀ لري۔
\begin{prooftree}
\AxiomC{$\Discharge{\lnot !A \lor !B}{1}$}
\DeduceC{$!A \lif !B$}
\DischargeRule{\Intro{\lif}}{1}
\UnaryInfC{$(\lnot !A \lor !B) \lif (!A \lif !B)$}
\end{prooftree}
له دې وروسته د دوام دوه لارې لرو۔ يا له لاندې څخه پورته کار روان
ساتلے او د~\Intro{\lif} بله کارونه لټولے شو، يا له پاسه لاندې کار
کولے او د~\Elim{\lor} قاعده کارولے شو۔ دويمه لار اخلو۔ فرضيه
$\lnot !A \lor !B$ د~\Elim{\lor} د تر ټولو کيڼې مقدمې په توګه
کاروو۔ د~\Elim{\lor} د سمې کارونې دپاره نورې دوه مقدمې بايد له
پايلې~$!A \lif !B$ سره يو شان وي، خو هره يوه په خپل وار له بلې
فرضیې، يعنې د~$\lnot !A \lor !B$ له دوو بېلندونو څخه له يوه،
رايستل کېدے شي۔ نو !!{derivation} به داسې وي:
\begin{prooftree}
\AxiomC{$\Discharge{\lnot !A \lor !B}{1}$}
\AxiomC{$\Discharge{\lnot !A}{2}$}
\DeduceC{$!A \lif !B$}
\AxiomC{$\Discharge{!B}{2}$}
\DeduceC{$!A \lif !B$}
\DischargeRule{\Elim{\lor}}{2}
\TrinaryInfC{$!A \lif !B$}
\DischargeRule{\Intro{\lif}}{1} 
\UnaryInfC{$(\lnot !A \lor !B) \lif (!A \lif !B)$}
\end{prooftree}

په ښي لوري دواړو څانګو کښې غواړو $!A \lif !B$~!!{derive} کړو؛
غوره لار ئې د~\Intro{\lif} کارول دي۔
\begin{prooftree}
\AxiomC{$\Discharge{\lnot !A \lor !B}{1}$}
\AxiomC{$\Discharge{\lnot !A}{2}, \Discharge{!A}{3}$}
\DeduceC{$!B$}
\DischargeRule{\Intro{\lif}}{3}
\UnaryInfC{$!A \lif !B$}
\AxiomC{$\Discharge{!B}{2}, \Discharge{!A}{4}$}
\DeduceC{$!B$}
\DischargeRule{\Intro{\lif}}{4}
\UnaryInfC{$!A \lif !B$}
\DischargeRule{\Elim{\lor}}{2}
\TrinaryInfC{$!A \lif !B$}
\DischargeRule{\Intro{\lif}}{1} 
\UnaryInfC{$(\lnot !A \lor !B) \lif (!A \lif !B)$}
\end{prooftree}

د !!{derivation} په دوو تشو برخو کښې له~$\lnot !A$ او~$!A$ څخه،
او له~$!A$ او~$!B$ څخه، د~$!B$ !!{derivation}s غواړو۔ لومړۍ برخه
اخلو۔ $\lnot !A$ او~$!A$ د~\Elim{\lnot} دوه مقدمې دي:
\begin{prooftree}
\AxiomC{$\Discharge{\lnot !A}{2}$}
\AxiomC{$\Discharge{!A}{3}$}
\RightLabel{\Elim{\lnot}}
\BinaryInfC{$\lfalse$}
\DeduceC{$!B$}
\end{prooftree}
د~\FalseInt په کارولو $!B$ د پايلې په توګه اخستلے او څانګه بشپړولے
شو۔
\begin{prooftree}
\AxiomC{$\Discharge{\lnot !A \lor !B}{1}$}
\AxiomC{$\Discharge{\lnot !A}{2}$}
\AxiomC{$\Discharge{!A}{3}$}
\RightLabel{\Elim{\lnot}}
\BinaryInfC{$\lfalse$}
\RightLabel{\FalseInt}
\UnaryInfC{$!B$}
\DischargeRule{\Intro{\lif}}{3}
\UnaryInfC{$!A \lif !B$}
\AxiomC{$\Discharge{!B}{2}, \Discharge{!A}{4}$}
\DeduceC{$!B$}
\DischargeRule{\Intro{\lif}}{4}
\UnaryInfC{$!A \lif !B$}
\DischargeRule{\Elim{\lor}}{2}
\TrinaryInfC{$!A \lif !B$}
\DischargeRule{\Intro{\lif}}{1} 
\UnaryInfC{$(\lnot !A \lor !B) \lif (!A \lif !B)$}
\end{prooftree}
\textbf{د ژباړې د نسخې سمون (OLND-003):} سرچينه په دې بشپړه ونه
کښې له نفي شوې مقدمې او مثبتې مقدمې څخه کاذبۍ ته استنتاج د کاذبۍ
د داخلولو په نښان ښيي؛ د مخکښې جزوي ونې او ټاکلې قاعدې سره سم دلته
د نفي ايستلو نښان راوړل شوے دے۔

اوس تر ټولو ښۍ څانګې ته ګورو۔ دلته بايد پوه شو چې د !!{derivation}
تعريف \emph{د فرضيو ايستل روا بولي} خو \emph{لازمي ئې نۀ بولي}۔ په
بله وينا، که له فرضيو~$!A$ او~$!B$ څخه د يوې په کارولو $!B$ راوباسو
او بله و نۀ کاروو، دا روا ده۔ له~$!B$ څخه د~$!B$~!!{derive} کول
څرګند دي: $!B$ په يوازې ځان داسې !!a{derivation} دے او هېڅ استنتاج
نۀ غواړي۔ نو فرضيه~$!A$ ساده ړنګولے شو۔
\begin{prooftree}
\AxiomC{$\Discharge{\lnot !A \lor !B}{1}$}
\AxiomC{$\Discharge{\lnot !A}{2}$}
\AxiomC{$\Discharge{!A}{3}$}
\RightLabel{\Elim{\lnot}}
\BinaryInfC{$\lfalse$}
\RightLabel{\FalseInt}
\UnaryInfC{$!B$}
\DischargeRule{\Intro{\lif}}{3}
\UnaryInfC{$!A \lif !B$}
\AxiomC{$\Discharge{!B}{2}$}
\RightLabel{\Intro{\lif}}
\UnaryInfC{$!A \lif !B$}
\DischargeRule{\Elim{\lor}}{2}
\TrinaryInfC{$!A \lif !B$}
\DischargeRule{\Intro{\lif}}{1} 
\UnaryInfC{$(\lnot !A \lor !B) \lif (!A \lif !B)$}
\end{prooftree}
پام وکړئ چې په بشپړ !!{derivation} کښې د تر ټولو ښي اړخ
\Intro{\lif} استنتاج په حقيقت کښې هېڅ فرضيه نۀ باسي۔
\end{ex}

\begin{ex}
تر اوسه مو د~\FalseCl{} قاعدې ته اړتيا نۀ ده پېښه شوې۔ دا ځکه
ځانګړې ده چې داسې فرضيه ايستل راکوي چې د قاعدې د پايلې فرعي
!!{formula} نۀ وي۔ دا له~\FalseInt{} قاعدې سره نژدې تړلې ده۔ په
حقيقت کښې \FalseInt{} د~\FalseCl{} قاعدې ځانګړے حالت دے---د
“شهودي منطق” په نامه يو منطق شته چې يوازې \FalseInt{} پکې روا دے۔
کله چې بله هېڅ لار کار نۀ کوي، \FalseCl{} وروستۍ وسيله ده۔ مثلاً
فرض کړئ غواړو $!A \lor \lnot !A$~!!{derive} کړو۔ زموږ عادي تګلاره
به دا وي چې د~$!A \lor \lnot !A$ د~$\Intro{\lor}$ په کارولو
!!{derive} کړو۔ خو دا
غواړي چې له هېڅ فرضیې پرته يا~$!A$ يا~$\lnot !A$~!!{derive} کړو،
او دا نۀ کېږي۔ \FalseCl{} مو ژغوري!
\begin{prooftree}
  \AxiomC{$\Discharge{\lnot(!A \lor \lnot !A)}{1}$}
  \DeduceC{$\lfalse$}
  \DischargeRule{\FalseCl}{1}
  \UnaryInfC{$!A \lor \lnot !A$}
\end{prooftree}
اوس له~$\lnot(!A \lor \lnot !A)$ څخه د~$\lfalse$ يو
!!a{derivation} لټوو۔ ځکه $\lfalse$ د~$\Elim{\lnot}$ پايله ده، نو
کېدے شي دا هڅه وکړو:
\begin{prooftree}
  \AxiomC{$\Discharge{\lnot(!A \lor \lnot !A)}{1}$}
  \DeduceC{$\lnot !A$}
  \AxiomC{$\Discharge{\lnot(!A \lor \lnot !A)}{1}$}
  \DeduceC{$!A$}
  \RightLabel{\Elim{\lnot}}
  \BinaryInfC{$\lfalse$}
  \DischargeRule{\FalseCl}{1}
  \UnaryInfC{$!A \lor \lnot !A$}
\end{prooftree}
د~$\lnot !A$ د !!a{derivation} موندلو تګلاره د~$\Intro{\lnot}$
کارونه غواړي:
\begin{prooftree}
  \AxiomC{$\Discharge{\lnot(!A \lor \lnot !A)}{1}, \Discharge{!A}{2}$}
  \DeduceC{$\lfalse$}
  \DischargeRule{\Intro{\lnot}}{2}
  \UnaryInfC{$\lnot !A$}
  \AxiomC{$\Discharge{\lnot(!A \lor \lnot !A)}{1}$}
  \DeduceC{$!A$}
  \RightLabel{\Elim{\lnot}}
  \BinaryInfC{$\lfalse$}
  \DischargeRule{\FalseCl}{1}
  \UnaryInfC{$!A \lor \lnot !A$}
\end{prooftree}
دلته $\lfalse$ په اسانۍ اخستلے شو: $\Elim{\lnot}$ پر
$\lnot(!A \lor \lnot !A)$ او~$!A \lor \lnot !A$ کاروو؛ دويمه
جمله زموږ له نوې فرضيې~$!A$ څخه د~$\Intro{\lor}$ په وسيله راځي:
\begin{prooftree}
  \AxiomC{$\Discharge{\lnot(!A \lor \lnot !A)}{1}$}
  \AxiomC{$\Discharge{!A}{2}$}
  \RightLabel{\Intro{\lor}}
  \UnaryInfC{$!A \lor \lnot !A$}
  \RightLabel{\Elim{\lnot}}
  \BinaryInfC{$\lfalse$}
  \DischargeRule{\Intro{\lnot}}{2}
  \UnaryInfC{$\lnot !A$}
  \AxiomC{$\Discharge{\lnot(!A \lor \lnot !A)}{1}$}
  \DeduceC{$!A$}
  \RightLabel{\Elim{\lnot}}
  \BinaryInfC{$\lfalse$}
  \DischargeRule{\FalseCl}{1}
  \UnaryInfC{$!A \lor \lnot !A$}
\end{prooftree}
په ښي لوري کښې همدا تګلاره کاروو، خو $!A$ د~\FalseCl په وسيله
اخلو:
\begin{prooftree}
  \AxiomC{$\Discharge{\lnot(!A \lor \lnot !A)}{1}$}
  \AxiomC{$\Discharge{!A}{2}$}
  \RightLabel{\Intro{\lor}}
  \UnaryInfC{$!A \lor \lnot !A$}
  \RightLabel{\Elim{\lnot}}
  \BinaryInfC{$\lfalse$}
  \DischargeRule{\Intro{\lnot}}{2}
  \UnaryInfC{$\lnot !A$}
  \AxiomC{$\Discharge{\lnot(!A \lor \lnot !A)}{1}$}
  \AxiomC{$\Discharge{\lnot !A}{3}$}
  \RightLabel{\Intro{\lor}}
  \UnaryInfC{$!A \lor \lnot !A$}
  \RightLabel{\Elim{\lnot}}
  \BinaryInfC{$\lfalse$}
  \DischargeRule{\FalseCl}{3}
  \UnaryInfC{$!A$}
  \RightLabel{\Elim{\lnot}}
  \BinaryInfC{$\lfalse$}
  \DischargeRule{\FalseCl}{1}
  \UnaryInfC{$!A \lor \lnot !A$}
\end{prooftree}
\end{ex}

\begin{prob}
داسې !!{derivation}s ورکړئ چې لاندې وښيي:
\begin{enumerate}
\item $!A \land (!B \land !C) \Proves (!A \land !B) \land !C$.
\item $!A \lor (!B \lor !C) \Proves (!A \lor !B) \lor !C$.
\item $!A \lif (!B \lif !C) \Proves !B \lif (!A \lif !C)$.
\item $!A \Proves \lnot\lnot !A$.
\end{enumerate}
\end{prob}

\begin{prob}
داسې !!{derivation}s ورکړئ چې لاندې وښيي:
\begin{enumerate}
\item $(!A \lor !B) \lif !C \Proves !A \lif !C$.
\item $(!A \lif !C) \land (!B \lif !C) \Proves (!A \lor !B) \lif !C$.
\item $\Proves \lnot(!A \land \lnot !A)$.
\item $!B \lif !A \Proves \lnot !A \lif \lnot !B$.
\item $\Proves (!A \lif \lnot !A) \lif \lnot !A$.
\item $\Proves \lnot(!A \lif !B) \lif \lnot !B$.
\item $!A \lif !C \Proves \lnot (!A \land \lnot !C)$.
\item $!A \land \lnot !C \Proves \lnot (!A \lif !C)$.
\item $!A \lor !B, \lnot !B \Proves !A$.
\item $\lnot !A \lor \lnot !B \Proves \lnot(!A \land !B)$.
\item $\Proves (\lnot !A \land \lnot !B) \lif\lnot(!A \lor !B)$.
\item $\Proves \lnot(!A \lor !B) \lif (\lnot !A \land \lnot !B)$.
\end{enumerate}
\end{prob}

\begin{prob}
داسې !!{derivation}s ورکړئ چې لاندې وښيي:
\begin{enumerate}
\item $\lnot(!A \lif !B) \Proves !A$.
\item $\lnot(!A \land !B) \Proves \lnot !A \lor \lnot !B$.
\item $!A \lif !B \Proves \lnot !A \lor !B$.
\item $\Proves \lnot \lnot !A \lif !A$.
\item $!A \lif !B, \lnot !A \lif !B \Proves !B$.
\item $(!A \land !B) \lif !C \Proves (!A \lif !C) \lor (!B \lif !C)$.
\item $(!A \lif !B) \lif !A \Proves !A$.
\item $\Proves (!A \lif !B) \lor (!B \lif !C)$.
\end{enumerate}
(دا ټول د~$\FalseCl$ قاعده غواړي۔)
\end{prob}

\end{document}
